\chapter*{Abstract}
\addcontentsline{toc}{chapter}{Abstract}

Large language models (LLMs) enable flexible interpretation and task decomposition, but linguistic plausibility does not establish that a multi-agent plan is executable, authorized, or complete. Conversely, symbolic artificial intelligence provides explicit constraints and verifiable state transitions but is less suited to interpreting underspecified natural-language goals. This thesis addresses the problem of coordinating heterogeneous autonomous agents without allowing an LLM proposal to substitute for an executable justification. Following a design-science methodology, it develops a protocol-independent framework that separates linguistic proposal from symbolic authorization. The framework represents agent capabilities, task dependencies, authority constraints, commitments, artifact contracts, and completion evidence explicitly. It returns either an authorized coordination plan or a structured infeasibility report. The framework is instantiated as a Python \texttt{CoordinationAgent} over \texttt{CoordinationSdk} middleware: the coordination agent owns planning and authorization, while the SDK owns communication, runtime adaptation, and normalization of remote, local, and optional Lingo-based linguistic agents. Evaluation uses unit and integration tests, seven deterministic coordination scenarios, an eleven-check Docker A2A system harness, and batched local LLM reference runs with \texttt{qwen/qwen3-1.7b} and \texttt{google/gemma-4-e2b}. The evidence supports the main boundary claim: language can propose structured candidates, but symbolic checks authorize or refuse execution and all dispatch proceeds through the SDK-managed communication surface. The work does not claim broad benchmark superiority over all possible baselines; its principal contribution is a precise and reproducible operational boundary for hybrid LLM--symbolic coordination.
